\documentclass[a4paper]{article}

% Español (México) con XeLaTeX: fontspec para fuentes Unicode, polyglossia
% para la división silábica y los nombres del documento en la variante mexicana.
\usepackage{fontspec}
\usepackage{polyglossia}
\setdefaultlanguage[variant=mexican]{spanish}
% Los nombres chinos van como «pinyin (original)»: xeCJK da a los caracteres
% Han su propia fuente; sin ella XeLaTeX los omite con un aviso.
% FandolSong no cubre algunos caracteres de nombres (U+73FA); WenQuanYi Zen Hei sí.
\usepackage[AutoFallBack=true]{xeCJK}
\setCJKmainfont{FandolSong-Regular.otf}
\setCJKfallbackfamilyfont{\CJKrmdefault}{WenQuanYi Zen Hei}
% polyglossia no traduce los nombres de listings.
\AtBeginDocument{\providecommand{\lstlistingname}{}\renewcommand{\lstlistingname}{Listado}%
  \providecommand{\lstlistlistingname}{}\renewcommand{\lstlistlistingname}{Índice de listados}}
\usepackage{amsmath,amssymb}
\usepackage{graphicx}
\usepackage[margin=2.5cm]{geometry}
\usepackage[most]{tcolorbox}
\usepackage{etoolbox}
\usepackage{subcaption}
\usepackage{listings}
\usepackage{hyperref}
\usepackage{booktabs}
\usepackage{float}

\newtcolorbox{knowledgebox}[1]{
    enhanced, colback=blue!5!white, colframe=blue!75!black, colbacktitle=blue!75!black,
    coltitle=white, fonttitle=\bfseries, title={#1}, boxrule=1pt, sharp corners
}

\newtcolorbox{importantbox}[1]{
    enhanced, colback=yellow!10!white, colframe=yellow!80!black, colbacktitle=yellow!80!black,
    coltitle=black, fonttitle=\bfseries, title={#1}, sharp corners
}

\newtcolorbox{warningbox}[1]{
    enhanced, colback=red!5!white, colframe=red!75!black, colbacktitle=red!75!black,
    coltitle=white, fonttitle=\bfseries, title={#1}, sharp corners
}

\lstset{
    basicstyle=\ttfamily\small,
    keywordstyle=\color{blue},
    stringstyle=\color{red!60!black},
    commentstyle=\color{green!60!black},
    breaklines=true,
    frame=single,
    numbers=left,
    numberstyle=\tiny\color{gray},
    captionpos=b,
    extendedchars=true
}

\newcommand{\notetitle}{CS224R Lecture 1: Introducción al aprendizaje por refuerzo profundo}
\newcommand{\noteauthors}{Elaboradas a partir de materiales públicos del curso}
\newcommand{\notedate}{Primavera de 2025}
\newcommand{\videochannel}{Stanford Online}
\newcommand{\videopublishdate}{2025-04}
\newcommand{\videoduration}{aproximadamente 80 minutos}
\newcommand{\videourl}{https://cs224r.stanford.edu/spring\_2025/}
\newcommand{\videocoverpath}{cover.jpg}
\newcommand{\repourl}{https://github.com/hqhq1025/ai-course-notes}


% Footer with repo info
\usepackage{fancyhdr}
\pagestyle{fancy}
\fancyhf{}
\fancyfoot[C]{\thepage}
\fancyfoot[R]{\footnotesize\color{gray}\href{\repourl}{AI Course Notes}}
\renewcommand{\headrulewidth}{0pt}
\renewcommand{\footrulewidth}{0pt}

\begin{document}

\begin{titlepage}
\centering
{\Large Notas del curso\par}
\vspace{1.2cm}
{\huge\bfseries \notetitle\par}
\vspace{0.8cm}
{\large \noteauthors\par}
\vspace{0.3cm}
{\large \notedate\par}
\vspace{1.1cm}

\ifdefempty{\videocoverpath}{
{\small Imagen de portada no proporcionada.\par}
}{
\includegraphics[width=0.82\textwidth,height=0.45\textheight,keepaspectratio]{\videocoverpath}\par
}

\vfill
\begin{tcolorbox}[width=0.9\textwidth, colback=black!2!white, colframe=black!60, sharp corners]
\textbf{Autor o canal del video}: \videochannel\par
\textbf{Fecha de publicación}: \videopublishdate\par
\textbf{Duración del video}: \videoduration\par
\textbf{Ponente}: Chelsea Finn\par
\textbf{Página del curso}: \href{https://cs224r.stanford.edu/spring_2025/}{cs224r.stanford.edu}\par
\textbf{Página del proyecto}: \href{\repourl}{\nolinkurl{\repourl}}\par
\end{tcolorbox}
\end{titlepage}

\tableofcontents
\newpage

%% --- Inicio del contenido --- %%

\section{Ubicación del curso: qué es el aprendizaje por refuerzo profundo}

La primera clase de CS224R parte de la pregunta más fundamental: qué es el aprendizaje por refuerzo profundo (Deep Reinforcement Learning), en qué se distingue del aprendizaje automático tradicional y por qué vale la pena estudiarlo.

\begin{importantbox}{Definición del aprendizaje por refuerzo profundo}
El aprendizaje por refuerzo profundo estudia los \textbf{problemas de decisión secuencial} (Sequential Decision-Making Problems): el sistema necesita, con base en un flujo de información, alternar continuamente entre "observar" y "actuar". El curso estudia también las \textbf{soluciones} a este tipo de problemas, entre ellas imitation learning, online/offline RL, model-free/model-based RL, multi-task/meta RL, así como RL for LLMs y RL for robots. Todos los métodos insisten en poder escalar a redes neuronales profundas.
\end{importantbox}

\subsection{Diferencia central con el aprendizaje supervisado}

En el aprendizaje supervisado contamos con datos etiquetados $\{(x_i, y_i)\}$, y el objetivo es aprender una función $f(x) \approx y$. Los puntos de datos son independientes e idénticamente distribuidos (i.i.d.), y se nos indica \textbf{directamente} cuál es la salida correcta.

En el aprendizaje por refuerzo, la situación es completamente distinta:

\begin{itemize}
    \item \textbf{Objetivo}: aprender una política de comportamiento $\pi(a \mid s)$ que mapee estados a acciones.
    \item \textbf{Retroalimentación}: no se indica directamente la respuesta, sino que se aprende de la experiencia mediante \textbf{retroalimentación indirecta} (como una señal de recompensa).
    \item \textbf{Los datos no son i.i.d.}: la acción $a$ del agente afecta la observación futura $s'$, de modo que la distribución de los datos depende de la política actual.
\end{itemize}

\begin{knowledgebox}{Las múltiples formas del comportamiento}
El "comportamiento" en el aprendizaje por refuerzo profundo puede incluir: el control de movimiento de un robot, la conversación de un chatbot, la estrategia de un juego, las decisiones de un vehículo autónomo, las operaciones de un Web agent, entre otros. El rasgo común de estos escenarios es que la salida (acción) del modelo produce consecuencias en el entorno, lo que a su vez afecta las entradas futuras.
\end{knowledgebox}

\subsection{Alcance y objetivos del curso}

Chelsea Finn señala explícitamente que este curso no puede abarcar todo el contenido del aprendizaje por refuerzo profundo. El curso se centra en:
\begin{itemize}
    \item Los \textbf{conceptos centrales} detrás de los métodos de RL profundo, de modo que los estudiantes puedan entender métodos más avanzados.
    \item La \textbf{implementación} de los algoritmos, para que los estudiantes puedan programarlos por sí mismos.
    \item El \textbf{control de robots y los modelos de lenguaje} como ejemplos principales, aunque la técnica tiene una aplicabilidad mucho más amplia.
\end{itemize}

\subsection{Resumen de la sección}

El aprendizaje por refuerzo profundo es la disciplina que trata de la decisión secuencial. En comparación con el aprendizaje supervisado, el RL enfrenta retos particulares: la retroalimentación es indirecta, los datos no son i.i.d. y las acciones tienen consecuencias. El objetivo central del curso es que los estudiantes puedan entender e implementar los métodos de RL profundo, tanto los existentes como los emergentes.
\section{Por qué estudiar el aprendizaje por refuerzo profundo}

\subsection{Más allá del aprendizaje supervisado}

Muchas aplicaciones de IA del mundo real no son simples mapeos de entrada a salida. El expositor da varios escenarios clave:

\begin{enumerate}
    \item \textbf{Las decisiones tienen consecuencias}: por ejemplo, después de que Spotify recomienda una canción, la siguiente recomendación debe evitar repetir la misma canción y debe mantener continuidad de estilo.
    \item \textbf{No se dispone de supervisión directa}: por ejemplo, en los asistentes de codificación (Cursor, Copilot), el usuario solo puede dar retroalimentación de "bueno" o "malo", en lugar de proporcionar directamente la salida correcta.
    \item \textbf{El objetivo no es diferenciable}: muchos objetivos reales (como la satisfacción del usuario) no se pueden optimizar directamente mediante el gradiente.
\end{enumerate}

\begin{figure}[H]
\centering
\includegraphics[width=0.9\textwidth]{slides-images/slide-18.jpg}
\caption{El curso usa la perspectiva de que "el despliegue cambia las observaciones futuras" para explicar por qué los problemas de decisión secuencial no pueden reducirse a un aprendizaje supervisado ordinario (Slide 18).}
\end{figure}

\subsection{Aplicación amplia en sistemas de IA de alto desempeño}

El expositor presentó una gran cantidad de aplicaciones reales del RL profundo:

\begin{itemize}
    \item \textbf{Robótica}: la marcha del robot cuadrúpedo de Unitree, la manipulación con $\pi_0$ de Physical Intelligence, Google Gemini Robotics.
    \item \textbf{Juegos}: la jugada 37 de AlphaGo mostró la capacidad del RL de descubrir estrategias nuevas que los humanos nunca habían imaginado.
    \item \textbf{Modelos de lenguaje grandes}: casi todos los LLM modernos usan alguna forma de RL para el post-training, sobre todo en el razonamiento avanzado.
    \item \textbf{Otros}: control de tráfico, seguimiento de prompt en modelos generativos de imágenes, diseño de chips (Google TPU).
\end{itemize}

\begin{figure}[H]
\centering
\includegraphics[width=0.88\textwidth]{slides-images/slide-19.jpg}
\caption{El RL profundo ya se ha convertido en un paradigma de entrenamiento importante para sistemas robóticos complejos, especialmente adecuado para tareas físicas que requieren control en lazo cerrado y adaptación en línea (Slide 19).}
\end{figure}

\begin{importantbox}{La capacidad de descubrimiento del RL}
El aprendizaje por refuerzo no solo imita el comportamiento ya presente en los datos, sino que también puede descubrir soluciones completamente nuevas mediante el "ensayo y error". La jugada 37 de AlphaGo es un caso clásico: fue una jugada que ningún jugador humano de go había hecho antes.
\end{importantbox}

\begin{figure}[H]
\centering
\includegraphics[width=0.88\textwidth]{slides-images/slide-22.jpg}
\caption{El curso también incluye explícitamente el post-training de LLM dentro del mapa de aplicaciones del RL profundo, lo que muestra que el RL se ha extendido de la robótica y los juegos a los sistemas de modelos de lenguaje (Slide 22).}
\end{figure}

\subsection{La naturaleza del aprendizaje}

Desde una perspectiva filosófica, aprender de la experiencia parece ser un componente fundamental de la inteligencia. Contar con algoritmos capaces de practicar por su cuenta y mejorar es esencial para construir sistemas verdaderamente inteligentes.

\subsection{Resumen de la sección}

Las razones para estudiar el RL profundo incluyen: (1) muchos problemas están fuera del alcance del aprendizaje supervisado; (2) los sistemas de IA de alto desempeño usan ampliamente el RL; (3) aprender de la experiencia es una capacidad fundamental de la inteligencia; (4) todavía hay una gran cantidad de problemas de investigación abiertos en el área.

\section{Modelado del comportamiento y fundamentos del aprendizaje por refuerzo}

\subsection{Representación de datos de la experiencia}

\begin{importantbox}{Términos clave}
\begin{itemize}
    \item \textbf{Estado} $s_t$: el estado del mundo en el instante $t$.
    \item \textbf{Observación} $o_t$: la información que el agente observa en el instante $t$ (se usa cuando hay pérdida de información).
    \item \textbf{Acción} $a_t$: la decisión en el instante $t$.
    \item \textbf{Trayectoria} $\tau = (s_1, a_1, s_2, a_2, \ldots, s_T, a_T)$: una secuencia completa de estados y acciones.
    \item \textbf{Función de recompensa} $r(s, a)$: mide qué tan buena es una pareja estado-acción.
\end{itemize}
\end{importantbox}

\subsection{Diferencia entre estado y observación}

Una propiedad clave es la \textbf{propiedad de Markov} (Markov Property): el siguiente estado depende únicamente del estado y la acción actuales, y no de la historia anterior:
$$p(s_{t+1} \mid s_t, a_t) \quad \text{independiente de } s_{t-1}$$

Pero en muchos escenarios reales (como los chatbots), el agente solo puede ver la \textbf{observación} y no el estado completo, por lo que hay que agregar memoria a la política: $\pi_\theta(a_t \mid o_{t-m}, \ldots, o_t)$.

\subsection{Representar el comportamiento con una red neuronal}

La política $\pi_\theta(a \mid s)$ es la función que mapea estados a una distribución sobre acciones. El proceso de ejecución del agente es el siguiente:
\begin{enumerate}
    \item Observa el estado $s_t$
    \item Muestrea una acción de la política $a_t \sim \pi_\theta(\cdot \mid s_t)$
    \item Observa el siguiente estado $s_{t+1} \sim p(\cdot \mid s_t, a_t)$
\end{enumerate}
La secuencia completa que se produce de esta manera se llama rollout de la política o episodio.

\subsection{El objetivo del aprendizaje por refuerzo}

El objetivo del aprendizaje por refuerzo es \textbf{maximizar la recompensa acumulada esperada}:
$$\max_\theta \; \mathbb{E}_{\tau \sim p_\theta(\tau)} \left[ \sum_{t} r(s_t, a_t) \right]$$

donde la distribución sobre trayectorias es:
$$p_\theta(\tau) = p(s_1) \prod_{t=1}^{T} \pi_\theta(a_t \mid s_t) \, p(s_{t+1} \mid s_t, a_t)$$

\begin{warningbox}{Por qué usar una esperanza}
La recompensa acumulada no es una cantidad determinista; tiene dos fuentes de aleatoriedad: (1) el entorno mismo es aleatorio; (2) la política puede ser aleatoria. Por eso lo que optimizamos es la recompensa acumulada \textbf{esperada}.
\end{warningbox}

\subsection{El sentido de las políticas aleatorias}

¿Por qué usar una política aleatoria en lugar de una política determinista? Por dos razones:
\begin{enumerate}
    \item \textbf{Exploración}: para aprender de la propia experiencia, hay que probar cosas distintas.
    \item \textbf{Modelar comportamiento aleatorio}: los datos existentes suelen mostrar comportamientos diversos, y se pueden aprovechar las herramientas de los modelos generativos.
\end{enumerate}

\subsection{La función de valor y la función Q}

Dos cantidades centrales para medir qué tan buena es una política:

\begin{itemize}
    \item \textbf{Función de valor} $V^\pi(s)$: la recompensa futura esperada al partir del estado $s$ y seguir la política $\pi$.
    \item \textbf{Función Q} $Q^\pi(s, a)$: la recompensa futura esperada al tomar la acción $a$ desde el estado $s$ y luego seguir la política $\pi$.
\end{itemize}

\subsection{Panorama de los tipos de algoritmos}

\begin{enumerate}
    \item \textbf{Imitation Learning}: imita una política capaz de obtener una recompensa alta.
    \item \textbf{Policy Gradients}: calcula directamente el gradiente de la función objetivo.
    \item \textbf{Actor-Critic}: estima el valor de la política actual y lo usa para mejorar la política.
    \item \textbf{Value-based}: estima el valor de la política óptima.
    \item \textbf{Model-based}: aprende un modelo de la dinámica del entorno, para usarlo en la planeación o en la mejora de la política.
\end{enumerate}

Los distintos algoritmos hacen distintas concesiones, y son adecuados para distintas condiciones supuestas (costo de recolección de datos, forma de la supervisión, estabilidad, características del espacio de acciones, etc.).

\subsection{Resumen de la sección}

Los elementos básicos del aprendizaje por refuerzo incluyen el estado, la acción, la trayectoria, la recompensa y la política. El objetivo es maximizar la recompensa acumulada esperada. La función de valor y la función Q son las herramientas centrales para evaluar qué tan buena es una política. Los distintos tipos de algoritmos tienen énfasis diferentes y son adecuados para escenarios diferentes.
\section{Síntesis y ampliación}

Esta clase estableció el marco general del aprendizaje por refuerzo profundo. Los puntos clave incluyen:

\begin{enumerate}
    \item El RL profundo estudia los problemas de decisión secuencial y sus soluciones, con una diferencia esencial respecto del aprendizaje supervisado.
    \item Las aplicaciones del RL ya son muy amplias: desde la robótica hasta los videojuegos y el post-training de los LLM.
    \item Marco de modelado básico: estado, acción, política, recompensa, trayectoria.
    \item El objetivo es maximizar la recompensa acumulada esperada $\mathbb{E}_{\tau \sim p_\theta(\tau)} \left[\sum_t r(s_t, a_t)\right]$.
    \item Las siguientes clases presentarán, en orden, imitation learning, policy gradients, actor-critic, Q-learning, offline RL, reward learning y RLHF.
\end{enumerate}

\subsection{Lecturas adicionales}
\begin{itemize}
    \item Sutton \& Barto, \textit{Reinforcement Learning: An Introduction} (2nd edition): libro de texto clásico de RL.
    \item CS234 (Stanford): un curso de RL más orientado a la teoría, complementario a CS224R.
    \item Sergey Levine, \textit{CS285 (Berkeley)}: otro excelente curso de RL profundo.
\end{itemize}

%% --- Fin del contenido --- %%

\end{document}
